\section{Related work}\label{sec:related}

\paragraph{Skill and luck in finance.} The distinction the benchmark operationalizes is old. Applying a false-discovery correction to three decades of mutual funds, \citet{barras2010false} find that roughly three quarters had no genuine alpha and that most apparent winners were false positives; \citet{kosowski2006stars} reach a more nuanced verdict by bootstrap, and \citet{fama2010luck} the same conclusion by a different route. The Sharpe ratio itself is due to \citet{sharpe1966mutual}, and \citet{sharpe1994sharpe} restates it as a differential return over a benchmark; this paper uses the zero-rate raw-return form (\cref{sec:stats}). \citet{lo2002sharpe} works out the sampling distribution of the Sharpe ratio and shows its square-root-of-time scaling fails off IID returns, which bears on the annualization that this paper's first finding concerns. \citet{bailey2014pseudo} show that a high backtest Sharpe is trivially achievable by trying enough configurations, \citet{harvey2015backtesting} turn the observation into a haircut that scales with the number of strategies tried, and \citet{harvey2016cross} apply the same discipline to the factor literature. The probabilistic Sharpe ratio of \citet{bailey2012sharpe} and the deflated Sharpe ratio of \citet{bailey2014deflated} are the instruments this benchmark uses; the Reality Check of \citet{white2000reality}, the superior-predictive-ability test of \citet{hansen2005spa} and the step-down procedure of \citet{romano2005stepwise} are the field-wide tests it reports. \citet{arnott2019protocol} set out a protocol for backtesting in the era of machine learning, and \cref{sec:experiments} follows its multiple-testing and reporting items: every Sharpe is reported beside its number of trials, its track length, and the cost profile it was earned under.

\paragraph{Forward and adaptive evaluation.} Committing to an evaluation before the data exists is not new, and the benchmark's intended forward protocol is positioned against that lineage. The M competitions run genuine forward windows on data that does not exist at submission time; M6 in particular scored forecasting and investment performance jointly over a year of live monthly submissions \citep{makridakis2025m6}. On the adaptive-overfitting side, \citet{dwork2015holdout} give a mechanism that keeps a holdout valid under repeated adaptive queries, and \citet{blum2015ladder} bound leaderboard overfitting from adaptive resubmission; both attack the same disease as deflation by controlling the query channel instead of correcting the statistic. At least one live trading arena for LLM agents exists \citep{deepfund2025}. What the chained board adds is independently checkable byte history: scoring rules and entries are bound before operator-declared deadlines, and a re-signed forgery is exposed by the next window's anchor (\cref{sec:attest}). The hash chain does not prove the deadline's wall time, data custody or prior non-observation. SharpeBench has no completed forward result and no multi-window operating history: its two empty pre-entry records were superseded when the scoring-provenance requirements changed, and its first live window is open but unresolved.

\paragraph{Benchmarks for financial agents.} \Cref{tab:related} positions the benchmark against the boards it is most often compared with, on the axes the principles in \cref{sec:principles} make load-bearing. \Cref{tab:related-inverse} then draws the same comparison on the axes the rivals hold: every rival that fields agents scores real LLM or learned agents, most present a far richer information environment, StockBench runs single-name equities on a genuinely post-cutoff window with a live board, and the Open FinLLM Leaderboard is hosted under neutral governance. SharpeBench's row in that table is empty, its empty cells are the limitations of \cref{sec:limitations}, and the two tables together state the actual trade: statistical and adversarial properties, each an executed check rather than a guarantee, were bought at the price of field breadth, and the program of \cref{sec:limitations} is to pay that price back without giving up the guarantees. Each mark in both tables records what the cited paper or its public board states, not an inference; the per-cell sources, including the deliberate blanks, are committed beside the evidence (\texttt{paper/evidence/table-provenance.md}). The cells marked ? are the ones still unchecked: costs for FinBen and FinRL-Meta, and single names for FinBen and FinRL-Meta. A cell marked $\circ$ records a partial form of the property: two boards repeat their runs and report the mean or the median, which is not a requirement that every run pass. The comparison therefore supports the statement that none of these boards states both deflation and $\passk$, not that each lacks every property its row leaves blank.

Each rival's ranking metric explains one of the axes. FinBen \citep{finben2024} ranks GPT-4 first on its trading task and reports its Sharpe beside buy-and-hold's with overlapping marginal 95 percent confidence intervals across ten stocks (\cref{sec:intro}). Overlapping marginal intervals are not a paired test: a Welch test on the same summary figures gives $p$ between $0.026$ and $0.050$ depending on how the half-widths are read, so the published numbers are borderline evidence of a difference rather than evidence that none exists, and they carry no correction for search. StockBench \citep{stockbench2025} evaluates on a single four-month post-cutoff window and ranks on cumulative return, drawdown and Sortino, with no deflation. QuantBench \citep{quantbench2025} names overfitting as an open problem and ranks pipeline performance with no luck-corrected metric, though its simulation does charge commissions and transaction fees. InvestorBench \citep{investorbench2024} evaluates thirteen LLM backbones as decision-making agents across stocks, crypto and ETFs, and ranks on return-based metrics with no deflation or reliability gate. FinRL-Meta \citep{finrlmeta2022} is scoped differently: it supplies hundreds of gym-style market environments and data pipelines for training RL agents, a benchmark of environments rather than a gated board of submitted agents, so most gate axes do not apply to it; it is in the tables because any serious agent field would plausibly be trained in it. The FINOS Open Financial LLM Leaderboard \citep{openfinllm} measures financial knowledge, scored on accuracy and F1, and has no trading-performance axis. The reproducibility gap across this field is quantified by the evidence map cited in \cref{sec:intro} \citep{llmtradingaudit2026}.

\paragraph{Evaluation practice beyond the compared boards.} Recent trading and finance papers outside \cref{tab:related} show the same gaps, and several report them about their own work. On repeated runs, Backtrader-Bench reports that a single run without tools overstated one model's ten-run mean by more than 12 points, 76.7 against 64.3 percent \citep[p.~4]{zhao2026backtrader}. Agora reports its single-seed numbers ``for parity with the AlphaGen reporting convention'' and has run its full system at one seed \citep[pp.~15 and 33]{li2026agora}. A NeurIPS 2025 trading paper answers its checklist question on error bars with ``Our evaluation environment is deterministic'' \citep[p.~17]{ma2025robust}, while its agents are trained with PPO \citep[p.~39]{ma2025robust}; \cref{sec:principles} states what a deterministic scorer does and does not guarantee. On costs, BacktestBench sets transaction costs and taxes to zero in its baseline experiments \citep[p.~3]{wang2026backtestbench}, and StockBench states that it does not model trading costs, slippage or liquidity limits \citep[arXiv version 2, p.~9]{stockbench2025}. On search, Agora lists multiple-testing correction among the analyses it defers \citep[p.~34]{li2026agora}. On leakage, two systems report finding their own. In an earlier version of AQuA's loop, a reviewer agent approved a feature normalized by the whole day's volume, and a manual audit found the leak after the held-out information coefficient came out far above that of comparable features \citep[pp.~21--22]{guo2026aqua}; in another pipeline, an off-by-one date slice put validation-year returns into the training labels and inflated test Sharpe by 0.3 \citep[p.~11]{du2026multifactor}. \Citet[p.~7]{gencay2026survives} plants a look-ahead oracle that reaches a DSR of $1.00$, the case \cref{sec:integrity} addresses by enforcing the look-ahead boundary at the agent interface rather than through a statistic. On judges, FinTradeBench reports a mean absolute error of 0.40 on a 1 to 5 scale between its LLM judge and human expert scores, beside a Spearman correlation of 0.07 that it attributes to a ceiling effect \citep[p.~25]{agrawal2026fintradebench}, and CryptoAnalystBench reports quadratic-weighted kappa between 0.35 and 0.42 between its judge and domain experts \citep[p.~8]{eswaran2026cryptoanalyst}. SharpeBench has no judge, so this kind of agreement check does not apply to its scores.

Re-scoring a rival field is infrastructure here, not a result. An import adapter is contributed that converts a rival board's published per-period return series into scoreable submissions, embedding the caveat that imported fields carry no audit trace and unknown trial counts understate deflation, retaining each column's own observed date axis so that two series missing different dates are not paired by position, and refusing a period repeated within one run, which would otherwise enter the mean, the variance, the PSR sample length and the bootstrap twice. The keyed scoring path applies the same refusal to a field that did not come through the adapter. No rival re-scoring has been executed and no demotion of any rival's leader is claimed: StockBench publishes environment inputs and summary statistics but no per-agent return series, so the experiment awaits either the authors' raw artifacts or a reproduction with their harness, and until then the adapter is infrastructure, not evidence.

\begin{table}[t]
\centering
\footnotesize
\caption{Positioning along the axes this benchmark adds. A \checkmark{} means the cited source states the property; a $\circ$ means it states a partial form, here repeated runs reported as a mean or a median rather than required to pass on every run; an empty cell means a search found no positive evidence; a cell marked ? is unchecked, with no exhaustive search of the full text, and is not a verified absence. Per-cell sources are in \texttt{paper/evidence/table-provenance.md}, except three cells read from the full text on 17 September 2026 and not yet in that sheet: StockBench (arXiv version 2) runs each agent three times with different seeds and reports the average (p.~5), InvestorBench (arXiv version 1) reports the trajectory with the median metrics of five repeated epochs (p.~6), and the full text of that InvestorBench version states no transaction cost, commission, fee or slippage. $\tau$-bench, a customer-service agent benchmark, appears because it is the construct source for the seed leg of the reliability gate. Read together with \cref{tab:related-inverse}: a comparison drawn only on one's own axes is advertising.}
\label{tab:related}
\setlength{\tabcolsep}{4pt}
\begin{tabular}{lcccccc}
\toprule
 & Deflates & $\passk$ & Process gate & Costs & Deterministic & Forward commit \\
\midrule
FinBen \citep{finben2024} & & & & ? & & \\
StockBench \citep{stockbench2025} & & $\circ$ & & & & \\
QuantBench \citep{quantbench2025} & & & & \checkmark & & \\
InvestorBench \citep{investorbench2024} & & $\circ$ & & & & \\
FinRL-Meta \citep{finrlmeta2022} & & & & ? & & \\
Open FinLLM \citep{openfinllm} & & & & & & \\
$\tau$-bench \citep{taubench2025} & & \checkmark & & & & \\
SharpeBench (this paper) & \checkmark & \checkmark & \checkmark & \checkmark & \checkmark & \checkmark \\
\bottomrule
\end{tabular}
\end{table}

\begin{table}[t]
\centering
\footnotesize
\caption{The same comparison on the axes the rival benchmarks hold. A mark means the cited source states the property; an empty cell means the provenance sheet records no positive evidence after a search; a cell marked ? is recorded there as unchecked, with no exhaustive search of the full text, and is not a verified absence (per-cell sources in \texttt{paper/evidence/table-provenance.md}). FinRL-Meta's near-empty row reflects its different scope, an environment library rather than an agent board. SharpeBench's empty cells are its limitations (\cref{sec:limitations}), and three of them, the missing LLM field, the thin information environment, and the absent live board, are the price paid for the guarantees in \cref{tab:related}.}
\label{tab:related-inverse}
\setlength{\tabcolsep}{4.5pt}
\begin{tabular}{lccccc}
\toprule
 & LLM agents & Rich info & Single names & Live board & Post-cutoff run \\
\midrule
FinBen \citep{finben2024} & \checkmark & \checkmark & ? & & \\
StockBench \citep{stockbench2025} & \checkmark & \checkmark & \checkmark & \checkmark & \checkmark \\
QuantBench \citep{quantbench2025} & \checkmark & & & & \\
InvestorBench \citep{investorbench2024} & \checkmark & & \checkmark & & \\
FinRL-Meta \citep{finrlmeta2022} & & & ? & & \\
Open FinLLM \citep{openfinllm} & \checkmark & \checkmark & & \checkmark & \\
$\tau$-bench \citep{taubench2025} & \checkmark & & & & \\
SharpeBench (this paper) & & & & & \\
\bottomrule
\end{tabular}
\end{table}

\paragraph{Evaluation as a scientific object.} The reliability-versus-capability split that motivates the seed leg of $\passk$ originates in \citet{taubench2025}, which asks whether an agent succeeds on all $k$ attempts rather than on at least one. It is now a theme across agent evaluation: \citet{miller2024errorbars} argues for error bars on every evaluation score, \citet{kapoor2024agents} for cost-controlled evaluation and adequate holdout sets, and \citet{agentreliability2026} treats reliability as an axis separate from capability and finds that gains in the latter have not bought the former. The surveys that shape this paper's integrity section, \citet{betterbench2024} and \citet{constructvalidity2025}, find that most benchmarks neither report the statistical significance of their rankings nor allow their results to be replicated, and recommend an operational construct-validity checklist, which \cref{app:checklist} answers. A 2026 survey of financial multi-agent evaluation \citep{finmultiagentsurvey2026} catalogues five recurring failure modes: look-ahead, survivorship, backtest overfitting, transaction-cost neglect and regime-shift blindness. SharpeBench addresses look-ahead through point-in-time observations, selection through deflation, transaction costs through the simulator and regime dependence through $\passk$. It does not close survivorship in the fixed universes, which \cref{sec:limitations} states.
